\documentclass[11pt,a4paper]{article}
\usepackage[margin=2.5cm]{geometry}
\usepackage{amsmath,amssymb,booktabs,graphicx,hyperref,natbib}
\usepackage[T1]{fontenc}
\usepackage{float}
\hypersetup{colorlinks=true,linkcolor=blue,citecolor=blue,urlcolor=blue}
\title{\textbf{Energy Shocks and the Green Adoption Channel in a Small Open Economy}}
\author{Guillaume Schwegler, Boris Chafwehe, Francesca Diluiso}
\date{\today}
\usepackage{mathtools}
\usepackage{xcolor}
% Marks a number that must be regenerated by re-running the code after the
% no-adoption counterfactual was changed to a common-steady-state comparison
% (see Section~\ref{sec:model}, "Constructing the no-adoption counterfactual").
% grep "NEEDSRERUN" to find every instance. The wrapped value is the OLD
% (pre-change) number, kept visible for reference, not a current result.
\newcommand{\NEEDSRERUN}[1]{\textcolor{red}{#1}}
\newcommand{\I}{\mathbb{I}}
\newcommand{\E}{\mathbb{E}}
\newcommand{\num}{\mathrm{num}}
\newcommand{\stst}{\mathrm{ss}}
\newcommand{\dsw}{D^{\mathrm{sw}}}
\newcommand{\DG}{D^{G}}
\newcommand{\PEB}{P^{E}_{B}}
\newcommand{\PEG}{P^{E}_{G}}
\newcommand{\pE}{p^{E}}
\newcommand{\coh}{\mathrm{coh}}

\begin{document}
\maketitle

\begin{abstract}
\noindent We add an extensive-margin brown/green durable adoption choice to the
small open economy HANK model of \citet{ARS2024}. Households choose discretely
whether to switch to a green durable that insulates them from fossil energy
prices, paying a resource cost. We study brown energy \emph{price} and
\emph{supply} shocks under the two fiscal responses studied in the recent
literature: a retail price cap \citep{Langot2026} and a Slutsky-compensating
transfer indexed to pre-crisis energy consumption \citep{Bayer2026}. The price
cap stabilises output but mechanically eliminates the green adoption margin,
while the transfer delivers comparable stabilisation and leaves adoption
intact. Shielding consumers from the relative price is therefore not neutral
for the energy transition. This is our contribution and, we think, its natural
boundary: the three closest papers---\citet{ARS2024}, \citet{Bayer2026} and
\citet{Langot2026}---all hold the household's energy technology fixed, and we
add the one margin they omit rather than the two-country or labour-market
machinery they contain. The transition loss is linear in the cap's intensity
under a price shock and convex under a supply shock. As a secondary
observation, the \emph{sign} of the cap's private-welfare effect flips with the
elasticity of energy supply, matching the sign each of \citet{Langot2026} and
\citet{Bayer2026} obtains in its own closure; we read this as locating one axis
of their contrast, not as nesting their models. The policy \emph{ordering} is
robust to calibration; the \emph{magnitude} of the adoption channel, and even
the sign of its output effect under a price shock, depend on the calibrated
switching cost and taste scale, which we report and discuss.
\end{abstract}

\section{Introduction}

Three recent papers study fiscal responses to the 2022 European energy crisis
in HANK frameworks. \citet{ARS2024} study a small open economy and stress the
negative externality when all importers subsidise simultaneously.
\citet{Bayer2026} use a two-country HANK with \emph{fixed} euro-area energy
supply and find subsidies are beggar-thy-neighbour while Slutsky transfers are
not. \citet{Langot2026} evaluate the French \emph{bouclier tarifaire} assuming
\emph{perfectly elastic} energy supply and find the price cap performs well.
All three treat the household's energy technology as fixed.

We relax exactly that margin, and only that margin. Households can adopt a
green durable that insulates them from the fossil price, at a resource cost.
The energy shock raises the return to adoption, so the crisis itself can
accelerate the transition---unless policy removes the price signal that drives
it. This is the contribution: an endogenous green-adoption margin, absent from
all three references. We deliberately do \emph{not} try to nest the two-country
spillovers of \citet{Bayer2026} or the labour-cost employment channel of
\citet{Langot2026}. A small open economy with an exogenous foreign block
contains neither, and competing on that machinery is not the point; our single
comparative advantage is the margin the literature leaves out, and we play only
that.

Because the price-versus-supply distinction turns out to matter for the
adoption margin, we do let the energy-supply elasticity vary: with
$\varepsilon^{E}=\infty$ the economy is a price taker (elastic supply, as in
Langot); with $\varepsilon^{E}$ finite the quantity is fixed and the world
price clears the market (as in Bayer). Both closures are reported below. This
is a modelling convenience for our own question, not a claim to reconcile the
two papers: the supply elasticity is one dimension along which they differ, and
not the one that carries their headline mechanisms (Section~\ref{sec:lit}).


\section{Model}\label{sec:model}

\subsection{Conventions and notation}
% ============================================================================

Time is quarterly. A subscript $t$ is suppressed where unambiguous; $x(-1)$ and
$x(+1)$ denote lag and (expected) lead. A bar denotes a steady-state value,
e.g.\ $\bar\pE$. All \emph{relative prices} carry the subscript ``\,$/P$\,''
where $P$ is the CPI, e.g.\ $\PEB\equiv \PEB{}_{,t}/P_t$; these are the objects
the CES demand system uses.

\paragraph{Unit of account.} The unit of account is the domestic good. The
household reads, solves in units of account, and returns assets converted to CPI
units by $A^{CPI}=A\,p^{\num}$, with
\begin{equation}
p^{\num}=p_{H/P},\qquad
1+r^{\num}=(1+r)\,\frac{p_{H/P}(-1)}{p_{H/P}},\qquad
atw^{\num}=\frac{atw}{p_{H/P}} .
\end{equation}
Objects contractually denominated in the CPI---fiscal transfers and the
switching cost $\psi_g$---are divided by $p^{\num}$ on the way into the budget.
Taking the domestic good rather than the CPI as unit of account makes the
resource constraint independent of how the CPI weights brown and green energy.

% ============================================================================
\subsection{Households}\label{sec:hh}
% ============================================================================

There are three permanent discount-factor types $i\in\{0,1,2\}$ of equal mass,
with $\beta_i$ (Section~\ref{sec:cal}); within a type, households face
idiosyncratic labour productivity and hold a discrete durable technology.

\subsubsection{Idiosyncratic productivity and the durable state}

Labour productivity $e$ follows a Rouwenhorst discretisation of an AR(1) in
logs, $\log e' = \rho_e \log e + \varepsilon$, with $n_e$ points, persistence
$\rho_e$ and innovation s.d.\ $\sigma_e$; transition matrix $\Pi$, mean-one grid.

Each household carries two binary durable indicators: the technology it
\emph{operates this period} $g\in\{B,G\}$ and the technology it \emph{held last
period} $g_-\in\{B,G\}$. We collapse the pair into a four-point index
\begin{equation}
d\in\{\,BB,\ BG,\ GB,\ GG\,\},\qquad
\text{(now, before)}=\{(B,B),(B,G),(G,B),(G,G)\}.
\end{equation}
Define the indicators
\begin{equation}
\mathrm{IS\_GREEN}(d)=\I[g=G]=(0,0,1,1),\qquad
\mathrm{PAYS\_SWITCH}(d)=\I[g=G,\,g_-=B]=(0,0,1,0).
\end{equation}
The switching cost $\psi_g$ is incurred exactly in the transition $g_-=B\to g=G$
(state $GB$). Brown is \emph{absorbing}; a green durable \emph{breaks down} to
brown at rate $\delta_g$. The breakdown Markov (applied in the \texttt{dep}
stage, rows = state this period, columns = state next period) is
\begin{equation}
M_{\mathrm{dep}}=
\begin{pmatrix}
1 & 0 & 0 & 0\\
1 & 0 & 0 & 0\\
\delta_g & 0 & 0 & 1-\delta_g\\
\delta_g & 0 & 0 & 1-\delta_g
\end{pmatrix}.
\end{equation}
On the equilibrium path $BG$ carries zero mass (green is lost only through
breakdown, recorded as $g_-=B$), so the effective state is three-point. In
steady state a permanent green share is sustained by a flow of switchers,
\begin{equation}
\dsw=\delta_g\,\DG. \label{eq:stockflow}
\end{equation}

Forward stage order within a period is
$\texttt{[dep, prod, durables, consav]}$: breakdown, productivity, the discrete
adoption choice, and the EGM consumption--savings step. The adoption stage uses
the discrete--continuous structure of \cite{DCEGM2017}.

\subsubsection{Preferences and the type-specific bundle price}

Flow felicity is CRRA over a consumption bundle $c$ with disutility of labour
handled at the union level (Section~\ref{sec:union}):
\begin{equation}
u(c)=\begin{cases}\log c & \text{eis}=1\\[2pt]
\dfrac{c^{1-1/\text{eis}}}{1-1/\text{eis}} & \text{eis}\neq 1.\end{cases}
\end{equation}
The bundle $c$ is a CES nest of energy against a home/foreign composite, with
ARS parameters $(\alpha_E,\eta_E)$. Because a green adopter buys energy at $\PEG$
and a brown household at $\PEB$, each durable type faces its own energy price and
hence its own \emph{exact} cost-of-living index:
\begin{align}
\pE_d &= \PEB + \mathrm{IS\_GREEN}(d)\,\bigl(\PEG-\PEB\bigr),\\
p_d &= \Bigl[\,1+\alpha_E\bigl((\pE_d)^{1-\eta_E}-(\PEB)^{1-\eta_E}\bigr)\Bigr]^{\frac{1}{1-\eta_E}}
\quad(\eta_E\neq 1),\qquad
p_d=\bigl(\pE_d/\PEB\bigr)^{\alpha_E}\ (\eta_E=1).
\label{eq:pd}
\end{align}
Equation~\eqref{eq:pd} uses the equilibrium CES identity
$\alpha_E(\PEB)^{1-\eta_E}+(1-\alpha_E)p_{HF/P}^{1-\eta_E}=1$
(the block \texttt{CESprices}, eq.~\eqref{eq:innernest}), so \emph{brown}
households have $p_d\equiv 1$ at every date (exact nesting). $p_d$ is
CPI-relative; the budget uses $p_d^{\num}\equiv p_d/p^{\num}$.

\subsubsection{Cash on hand}\label{sec:coh}

Let $a$ be beginning-of-period assets (in units of account, on a grid with lower
bound $a\ge \underline a$, $\underline a=0$). Cash on hand of a household in
durable state $d$ with productivity $e$ is
\begin{equation}
\coh(d,e,a)=(1+r^{\num})\,a+atw^{\num}\,e
+\frac{T^{\text{fisc}}}{p^{\num}}
-\frac{\psi_g}{p^{\num}}\,\mathrm{PAYS\_SWITCH}(d),
\label{eq:coh}
\end{equation}
with the (CPI-denominated) transfer
\begin{equation}
T^{\text{fisc}} = \epsilon^{T}+\iota\,(\pE-\bar\pE)\,\bar c^{B}_{E,i}(e,a)
\qquad\text{(untargeted lump sum $\epsilon^{T}$ + Slutsky transfer).}
\end{equation}
Here $\iota=\texttt{insE}$ scales the Slutsky transfer, $\epsilon^{T}=\texttt{epsT}$
is an untargeted lump sum, and $\bar c^{B}_{E,i}(e,a)$ is household $i$'s
steady-state \emph{brown} energy demand, collapsed over the durable axis and
broadcast identically across durable states (so a switcher keeps the same
transfer; the switching margin is undistorted). Both instruments are zero in the
baseline. The switching cost enters cash on hand only in state $GB$.

\subsubsection{Discrete adoption choice}\label{sec:logit}

In the \texttt{durables} stage the household chooses its current technology $g$
given the technology it holds (post-breakdown), by a logit over continuation
values with taste scale $\sigma_\varepsilon=\texttt{taste\_shock}$. Feasible
transitions carry the flow payoff
\begin{equation}
f(d\mid d_-)=
\begin{cases}
0 & \text{stay/renew brown: } (BB\mid BB),(BB\mid BG),\\
-\,\texttt{green\_block}\ \ (\text{masked to }-\infty\text{ if }\coh(GB)\le 0)
& \text{switch to green: } (GB\mid BB),(GB\mid BG),\\
0 & \text{stay/renew green: } (GG\mid GB),(GG\mid GG),\\
-\infty & \text{otherwise (infeasible).}
\end{cases}
\end{equation}
The switching cost itself is already in $\coh(GB)$ through
eq.~\eqref{eq:coh}. \texttt{green\_block} is a counterfactual wedge, $0$ in the
baseline (margin open) and $+\infty$ (large) in the ``no-adoption'' variant
(margin shut, used for the adoption-channel decomposition). Given post-choice
values $V(d)$, the choice probabilities and the expected (pre-shock) value are
the standard logit / log-sum-exp:
\begin{equation}
\Pr(d\mid d_-)=\frac{\exp\!\bigl([f(d\mid d_-)+V(d)]/\sigma_\varepsilon\bigr)}
{\sum_{d'}\exp\!\bigl([f(d'\mid d_-)+V(d')]/\sigma_\varepsilon\bigr)},
\qquad
\tilde V(d_-)=\sigma_\varepsilon\log\!\sum_{d'}\exp\!\Bigl(\tfrac{f(d'\mid d_-)+V(d')}{\sigma_\varepsilon}\Bigr).
\end{equation}

\subsubsection{Consumption--savings (EGM)}

Conditional on the durable state, the household solves a standard
consumption--savings problem with the type-specific bundle price $p_d^{\num}$.
The budget, Euler and envelope are
\begin{align}
p_d^{\num}\,c + a' &= \coh(d), & a'&\ge \underline a, \label{eq:budget}\\
u'(c) &= \beta_i\,p_d^{\num}\,\E\,V_a(a',e',d') & &\text{(Euler)},\\
V_a &= (1+r^{\num})\,u'(c)/p_d^{\num} & &\text{(envelope)},\\
V &= u(c)+\beta_i\,\E\,V(a',e',d') & &\text{(value)},
\end{align}
solved by the endogenous grid method, with $u'(c)=c^{-1/\text{eis}}$.

\subsubsection{Household outputs (hetoutputs)}\label{sec:hetout}

Given the policy $c(d,e,a)$, the household's \emph{exact} CES demand split by
durable type is
\begin{align}
c_E(d) &= \alpha_E\bigl(\pE_d/p_d\bigr)^{-\eta_E} c, &
c_{HF}(d) &= (1-\alpha_E)\bigl(p_{HF/P}/p_d\bigr)^{-\eta_E} c, \label{eq:demand}\\
c_E^{B}(d) &= c_E(d)\,\bigl(1-\mathrm{IS\_GREEN}(d)\bigr), &
c_E^{G}(d) &= c_E(d)\,\mathrm{IS\_GREEN}(d).
\end{align}
(Prices in \eqref{eq:demand} are CPI-relative; the denominator is $p_d$, not
$p_d^{\num}$.) The block also returns $d_G(d)=\mathrm{IS\_GREEN}(d)$,
$d_{sw}(d)=\mathrm{PAYS\_SWITCH}(d)$, the MPC, and flow utility $u(c)$.

% ============================================================================
\subsection{Aggregation}\label{sec:agg}
% ============================================================================

For each type $i$, SSJ aggregates the hetoutputs over the stationary
distribution $\Lambda_i$ into $C_i,A_i,\mathrm{MPC}_i,\DG_i,\dsw_i,
C^{B}_{E,i},C^{G}_{E,i},C_{HF,i},U_i$, and the model averages equally across the
three types:
\begin{equation}
X=\tfrac13\textstyle\sum_i X_i,\qquad
X\in\{C,A,\mathrm{MPC},\DG,\dsw,C^{B}_{E},C^{G}_{E},C_{HF},\ldots\},
\qquad \beta=\tfrac13\textstyle\sum_i\beta_i .
\end{equation}
Assets in units of account are converted to CPI units, $A^{CPI}=A\,p^{\num}$.
Aggregate demands for the home good, the foreign good and energy are then built
from the household aggregates (block \texttt{hh\_outputs\_dur}), applying on top
only the inner home/foreign nest, common to all households:
\begin{equation}
c_E = C^{B}_{E}+C^{G}_{E},\qquad
c_H = (1-\alpha_F)\,p_{H/HF}^{-\eta}\,C_{HF},\qquad
c_F = \alpha_F\,p_{F/HF}^{-\eta}\,C_{HF}. \label{eq:cHcF}
\end{equation}
Using $\sum_d c_E(d)$ rather than the CES demand at the aggregate index avoids a
Jensen bias of order $0.5\%$.

% ============================================================================
\subsection{Firms and domestic production}\label{sec:firms}
% ============================================================================

A competitive final domestic good is produced from labour. With productivity $Z$
and a constant gross markup $\mu_{\stst}$,
\begin{equation}
y = Z\,n,\qquad
\mathrm{gdp}=Z\,n,\qquad
\text{btw}_n = w\,n,\qquad
atw_n = (1-\tau^{Y})\,w\,n,
\end{equation}
where $y$ is home-good output, $n$ hours, $w$ the real wage. Monopolistic
profits are capitalised into the mutual fund:
\begin{equation}
\mathrm{div}=(\mu_{\stst}-1)\,w\,n,\qquad
J=\mathrm{div}+\frac{J(+1)}{1+r^{\text{ante}}},\qquad
j=\frac{J(+1)}{1+r^{\text{ante}}}.
\end{equation}

\subsubsection{Wage setting (sticky wages)}\label{sec:union}

A union sets the nominal wage subject to Calvo frictions $\theta_w$, giving a
wage Phillips curve. With $\kappa_w=(1-\theta_w)(1-\beta\theta_w)/\theta_w$,
labour-disutility shifter $\varphi$ (\texttt{vphi}), Frisch elasticity
$\texttt{frisch}$, and a real-wage-rigidity term controlled by $w_{BG}$,
\begin{align}
\pi^{w}(1+\pi^{w})
&=\beta\,\pi^{w}(+1)\bigl(1+\pi^{w}(+1)\bigr)
+\frac{\kappa_w}{\varphi\,n^{1/\text{frisch}}}\;\Xi,\\
\Xi &= \varphi\,n^{1/\text{frisch}}
-\frac{atw\,C^{-1/\text{eis}}}{\mu_{\stst}}
- w_{BG}\,\underbrace{\frac{(w-\bar w)\,w\,\bar C^{-1/\text{eis}}\,\bar n}{\mu_{\stst}\,n\,\bar w}}_{\text{real-wage rigidity}},
\end{align}
and the nominal and real wage evolve as $W=(1+\pi^{w})W(-1)$, $w=W/P$. Here
$atw=atw_n/n$ is the after-tax wage per hour.

% ============================================================================
\subsection{Prices}\label{sec:prices}
% ============================================================================

All retail prices are relative to the CPI $P$; $Q$ is the real exchange rate
(price of the foreign basket in home units), $P^{*}_{F},P^{E*}$ are exogenous
world prices.

\paragraph{Domestic good.}
\begin{equation}
P_H=\frac{\mu_{\stst}}{Z}\,W,\qquad
p_{H/P}=P_H/P,\qquad
\pi^{H}=P_H/P_H(-1)-1,\qquad w=W/P.
\end{equation}

\paragraph{Imported foreign good (Calvo $\theta_F$).} With
$\kappa_F=(1-\theta_F)(1-\beta_F\theta_F)/\theta_F$, $\beta_F=1/(1+\bar r^{\text{ante}})$,
\begin{equation}
\kappa_F\Bigl(\frac{Q P^{*}_{F}}{p_{F/P}}-1\Bigr)
+\beta_F\,\pi^{F}(+1)\bigl(1+\pi^{F}(+1)\bigr)
-\pi^{F}(1+\pi^{F})=0,\qquad p_{F/P}=P_F/P .
\end{equation}

\paragraph{Energy retail price and the price cap (Calvo $\theta_E$).} The
wholesale/retail brown energy price $\pE=P_E/P$ satisfies a NKPC, and the
\emph{household} brown price applies the cap:
\begin{align}
\kappa_E\Bigl(\frac{Q P^{E*}}{\pE}-1\Bigr)
+\beta_E\,\pi^{E}(+1)\bigl(1+\pi^{E}(+1)\bigr)
-\pi^{E}(1+\pi^{E})&=0,
& \kappa_E&=(1-\theta_E)(1-\beta_E\theta_E)/\theta_E,\\
\PEB &= (1-\tau^{E})\,\pE+\tau^{E}\,\bar\pE,
&& \tau^{E}=1\Rightarrow \PEB\ \text{pinned at }\bar\pE .
\label{eq:cap}
\end{align}

\paragraph{Green energy price.} With green/brown ratio $\varrho=\texttt{pE\_g\_ratio}$,
the green price is fixed relative to the pre-crisis brown price:
\begin{equation}
\PEG=\varrho\,\bar\pE . \label{eq:pEG}
\end{equation}

\paragraph{CPI (brown-anchored; ``Option~C'').} The energy leg of the CPI is
anchored on the brown price alone. With $\alpha=\alpha_E+(1-\alpha_E)\alpha_F$,
the inner (energy vs home/foreign) and outer (home vs foreign) CES identities
that define $p_{HF/P}$ and $p_{H/HF}$ are
\begin{align}
(1-\alpha_E)\,p_{HF/P}^{1-\eta_E}+\alpha_E\,(\PEB)^{1-\eta_E}&=1, \label{eq:innernest}\\
(1-\alpha_F)\,p_{H/HF}^{1-\eta}+\alpha_F\,p_{F/HF}^{1-\eta}&=1, \label{eq:outernest}
\end{align}
with $p_{F/HF}=p_{F/P}/p_{HF/P}$, $p_{H/P}=p_{H/HF}\,p_{HF/P}$, $p_H^{*}=p_{H/P}/Q$.
The gap between this measured CPI and the share-weighted (``true'') deflator is
computed \emph{ex post}, outside the model:
\begin{equation}
\phi^{1-\eta_E}=1+\alpha_E\,\DG\bigl[(\PEG)^{1-\eta_E}-(\PEB)^{1-\eta_E}\bigr].
\end{equation}
Under the domestic-good numeraire this choice affects only the measured deflator
and the monetary rule, not the resource constraint.

\paragraph{Aggregate price level and inflation.}
\begin{equation}
P=(1+\pi)P(-1),\qquad
\pi^{\text{ann}}=(1+\pi)^4-1,\ \text{etc.}
\end{equation}

% ============================================================================
\subsection{Open economy}\label{sec:open}
% ============================================================================

\paragraph{Foreign demand for the home good.} With foreign demand shifter
$\alpha^{*}$, trade elasticity $\gamma$, and world demand $C^{*}$,
\begin{equation}
c_H^{*}=\alpha^{*}\,\bigl(p_H^{*}\bigr)^{-\gamma}\,C^{*}.
\end{equation}

\paragraph{Uncovered interest parity.} The real exchange rate solves
\begin{equation}
\frac{Q}{Q(+1)}\,(1+r^{\text{ante}})=1+r^{*}. \label{eq:uip}
\end{equation}

\paragraph{Importer margins (foreign-owned).} Retail importers of the foreign
good and of energy are \emph{owned by foreigners} (ARS, \S4.1), so their margins
over the world price accrue abroad and are \emph{not} capitalised into domestic
wealth. They enter the model only by valuing the import bill at the retail (not
world) price in the balance of payments below:
\begin{equation}
D_F = \bigl(p_{F/P}-Q P^{*}_{F}\bigr)c_F,\qquad
D_E = \bigl(\pE-Q P^{E*}\bigr)c_E .
\end{equation}
Both margins are zero in steady state ($p_{F/P}=Q P^{*}_{F}$, $\pE=Q P^{E*}$), so
this convention leaves the steady state unchanged.

\paragraph{World energy supply (IEA block).} World supply is a fixed quantity
$E^{\text{shock}}$ plus intertemporal stock arbitrage with cost $\Gamma_{\text{arb}}$:
\begin{equation}
E_{\text{stock}} = \frac{P^{E*}(+1)/(1+r^{*})-P^{E*}}{\Gamma_{\text{arb}}},\qquad
E^{\text{supply}} = E^{\text{shock}}+\bigl(E_{\text{stock}}(-1)-E_{\text{stock}}\bigr).
\end{equation}
In the baseline the home economy owns \emph{none} of the world energy endowment
($\zeta^{E}=0$: a pure importer), so the associated rents are not capitalised
into domestic wealth. The endowment-ownership channel ($\zeta^{E}>0$, ARS
\S3.4/\S4.1) is retained in the code and used only as a robustness/energy-independence
experiment.

\paragraph{Fund return (revaluation).} The realised (ex post) return on the
fund's claims defines $r$:
\begin{equation}
1+r=\frac{J}{j(-1)}.
\end{equation}
Importers are foreign-owned and the energy endowment is not held domestically
($\zeta^{E}=0$), so the only domestically held risky claim is domestic firms:
$A^{\text{dom}}=j$, with return $r^{\text{dom}}=r$; $A^{\text{dom}}$ enters only
the current-account revaluation term below.

\paragraph{Balance of payments.}\label{sec:bop}
The green margin enters the import booking through two terms,
\begin{equation}
\text{imports}_{\text{dur}}=\psi_g\,\dsw
\quad\text{(switching cost booked as import)},\qquad
\text{energy\_gap}=(\PEB-\PEG)\,C^{G}_{E}
\quad\text{(adopters' import saving)}.
\end{equation}
Exports, imports and net exports are
\begin{align}
\text{exports}&=Q\,p_H^{*}c_H^{*},\\
\text{imports}&=p_{F/P}\,c_F+\pE\,c_E+\text{imports}_{\text{dur}}-\text{energy\_gap},\\
nx&=Q\,p_H^{*}c_H^{*}-p_{F/P}\,c_F-\pE\,c_E-\text{imports}_{\text{dur}}+\text{energy\_gap}.
\end{align}
Net foreign assets accumulate with a revaluation term:
\begin{align}
\text{reval} &= (r-r^{\text{ante}}(-1))A^{CPI}(-1)-(r^{\text{dom}}-r^{\text{ante}}(-1))A^{\text{dom}}(-1),\\
\mathrm{nfa} &= nx+\text{reval}+(1+r^{\text{ante}}(-1))\,\mathrm{nfa}(-1),
\qquad nx/\mathrm{gdp}=nx/y. \label{eq:nfa}
\end{align}

% ============================================================================
\subsection{Policy}\label{sec:policy}
% ============================================================================

\subsubsection{Monetary policy}

The nominal rate rule and the Fisher ex-ante real rate are
\begin{equation}
1+i = (1+\bar r^{*})\,(1+\phi_\pi\pi)\,(1+\phi_{\pi e}\pi(+1))+\text{ishock},
\qquad
r^{\text{ante}} = \frac{1+i}{1+\pi(+1)}-1.
\end{equation}
The baseline is the ARS \emph{constant real rate} rule
$(\phi_\pi,\phi_{\pi e})=(0,1)\Rightarrow r^{\text{ante}}=\bar r^{*}$; a standard
Taylor rule sets $(\phi_\pi,\phi_{\pi e})=(1.5,0)$. \texttt{ishock} is a monetary
policy shock (zero unless a monetary experiment is run).

\subsubsection{Fiscal policy}\label{sec:fiscal}

The government runs the two energy-crisis instruments and an untargeted lump sum,
all debt financed, with debt stabilised through the distortionary labour tax:
\begin{align}
\text{Subsidy} &= \tau^{E}\,(\pE-\bar\pE)\,c_E
&& \text{price cap (cost scales with actual use),}\\
\text{Ttargeted} &= \iota\,(\pE-\bar\pE)\,\bar C^{B}_{E}
&& \text{Slutsky transfer (cost scales with pre-crisis brown use),}\\
\text{Tuntargeted} &= \epsilon^{T}, &&\\[2pt]
\text{spending} &= \text{Tuntargeted}+\text{Ttargeted}+\text{Subsidy},\qquad
\text{taxation}=\tau^{Y}\,w\,n. &&
\end{align}
The government budget constraint and the debt-stabilising tax rule are
\begin{align}
B &= (1+r^{\text{ante}}(-1))\,B(-1)+\text{spending}-\text{taxation},
\label{eq:gbc}\\
\tau^{Y} &= \psi_B\bigl(B(-1)-\bar B\bigr),\qquad \psi_B=\texttt{psiB}.
\label{eq:taurule}
\end{align}
In the baseline $\bar B=0$ (no steady-state government debt): crisis outlays are
deficit-financed and repaid by raising the labour tax in proportion to the debt's
deviation from zero. Government debt $B$ is a genuine transition unknown (it moves
in the dynamics) and enters household wealth through asset-market clearing,
eq.~\eqref{eq:assets}.

% ============================================================================
\subsection{Market clearing and equilibrium}\label{sec:clearing}
% ============================================================================

The equilibrium residuals are
\begin{align}
\text{goods:}&\quad c_H+c_H^{*}-y=0, \label{eq:goods}\\
\text{assets (Walras):}&\quad A^{CPI}-\mathrm{nfa}-j-B=0, \label{eq:assets}\\
\text{energy:}&\quad
\begin{cases}
P^{E*}-P^{E*,\text{shock}}=0 & \varepsilon^{E}=\infty\ \text{(price taker; price shock)},\\
c_E-E^{\text{supply}}=0 & \varepsilon^{E}<\infty\ \text{(fixed quantity; supply shock).}
\end{cases}\label{eq:energy}
\end{align}
Equation~\eqref{eq:assets} is a Walras'-law residual (not itself a solved
target); it is monitored to machine precision as a consistency check.

\paragraph{Transition unknowns/targets.} The linear (SSJ) transition solves for
$\{y,\pi,r,\tau^{Y},P^{E*},w\}$ against
$\{\text{goods},\ \pi\text{ residual},\ r\text{ residual},\ \tau^{Y}\text{ residual},
\ \text{energy},\ w\text{ residual}\}$.

% ============================================================================
\subsection{Steady state and shocks}\label{sec:ss}
% ============================================================================

\subsubsection{Openness normalisations}

With $\zeta^{E}=0$ (pure importer) the ARS home-endowment rescaling is inactive,
so the normalisations reduce to
\begin{equation}
\alpha_F=\frac{\alpha-\alpha_E}{1-\alpha_E},\qquad
\gamma=\frac{\chi_{\text{tgt}}-(1-\alpha)(1-\alpha_F)\eta_E}{(1-\alpha)\alpha_F+\alpha},\quad \eta=\gamma,\ \ \chi_{\text{tgt}}=0.3,
\end{equation}
\begin{equation}
y=1,\qquad Z=\mu_{\stst}=1.03,\qquad \alpha^{*}=\alpha,\qquad E^{\text{shock}}=\alpha_E .
\end{equation}
(For the robustness runs with $\zeta^{E}>0$ the endowment rescaling
$y=1-\zeta^{E}\alpha_E$, $\mu_{\stst}\leftarrow\mu_{\stst}(1-\zeta^{E}\alpha_E)$,
$\alpha^{*}=\alpha-\zeta^{E}\alpha_E$ applies.)

\subsubsection{Steady-state unknowns/targets}

The steady state solves for $\{\varphi,\ \beta_{\max},\ y,\ \psi_g\}$ against
$\{\pi^{w}\text{ res},\ \mathrm{nfa}\text{ res},\ \text{goods},\ \DG-0.05\}$: $\psi_g$
is calibrated to the level target $\DG_{\stst}=0.05$ and $\beta_{\max}$ closes net
foreign assets. In the ``no-adoption'' counterfactual ($\texttt{green\_block}$
large) $\DG\equiv 0$ for every $\psi_g$, so $\psi_g$ is not identifiable and is
carried over, fixed, from the matching adoption steady state.

\subsubsection{Shocks}

\begin{align}
\text{Brown-price shock (requires }\varepsilon^{E}=\infty):&\quad
P^{E*,\text{shock}}_t=\text{size}\cdot\rho^{t},\quad \rho=2^{-1/\text{half-life}},\\
\text{Brown-supply shock (requires }\varepsilon^{E}<\infty):&\quad
E^{\text{shock}}_t=\bar E^{\text{shock}}-\text{drop}\cdot\bar E^{\text{shock}}\ \text{for }t<\text{quarters},\ \text{else }\bar E^{\text{shock}} .
\end{align}
The baseline price shock is $\text{size}=1$ (world price doubles on impact),
half-life $16$; the baseline supply shock is $-10\%$ for $6$ quarters.

\section{Calibration}\label{sec:cal}

\begin{table}[H]\centering
\small
\begin{tabular}{lll}
\toprule
Parameter & Value & Source / target \\
\midrule
$\alpha_E$ energy expenditure share & 0.04 & \citet{ARS2024} \\
$\eta_E$ energy substitution & 0.10 & \citet{ARS2024} \\
$r$, $\mu_{ss}$ & 0.01, 1.03 & \citet{ARS2024} \\
$\beta$ heterogeneity (3 types) & spread 0.06 & fixed; implied MPC $\approx0.30$ \\
\midrule
$\delta_g$ green breakdown rate & 0.05 & assumption (5\,yr life) \\
$\varrho=P^{E}_{G}/P^{E}_{B}$ steady-state gap & 0.80 & heat-pump / EV delivered-cost ratio \\
$\psi_g$ switching cost & 0.538 & targets $D^{G}_{ss}=0.05$ \\
$\sigma_{\varepsilon}$ taste scale & 0.05 & calibrated jointly with $\psi_g$ \\
$\rho_g$ pass-through to green price & 0 & upper bound on the channel \\
\bottomrule
\end{tabular}
\caption{Calibration. The steady state is identical across all experiments.
$\psi_g$ and $\sigma_\varepsilon$ are calibrated jointly to the level target
$D^{G}_{ss}=0.05$; $\sigma_\varepsilon$ is not separately pinned down by an
adoption-elasticity moment (Section~\ref{sec:limits}). The gross markup carries
ARS's home-endowment normalisation $\mu_{ss}\!\leftarrow\!\mu_{ss}(1-\zeta^E\alpha_E)$,
so the value entering firm profits is $\approx1.016$; $1.03$ is the
pre-normalisation primitive.}
\end{table}

\paragraph{Robustness to the green--brown cost ratio $\varrho$.} The one
durable-side parameter with a direct empirical counterpart is the steady-state
green/brown cost ratio $\varrho$. We anchor it at $0.80$ from the relative cost
of delivered energy service under adoption: a heat pump with a coefficient of
performance near three delivers heat from roughly a third of the fuel-equivalent
input, which at 2019--21 electricity-to-gas price ratios places the delivered
green cost around $0.7$--$0.8$ of brown; the per-kilometre cost of an electric
against a comparable internal-combustion vehicle is in the same range.
Table~\ref{tab:rho} re-solves the model across $\varrho\in\{0.6,0.7,0.8\}$,
re-calibrating $\psi_g$ each time to hold $D^{G}_{ss}=0.05$. The \emph{ordering}
is invariant: the cap collapses the peak green-share response to essentially
zero at every $\varrho$, the transfer leaves it at its no-policy value, the cap
cushions cumulative output and the transfer boosts it. Only magnitudes move, and
modestly (the no-policy peak green share ranges over $\NEEDSRERUN{8.4\text{--}9.2}$ pp). The
limiting case $\varrho=1$---green no cheaper than brown---is degenerate for the
level target: no positive switching cost can sustain a $5\%$ green steady state
when adoption confers no cost advantage, so adoption collapses, which is exactly
the sanity check one wants.

\begin{table}[H]\centering\small
\textcolor{red}{\bfseries\textdagger\ NEEDSRERUN --- this entire table is
hardcoded (not \texttt{\textbackslash input}-generated) and its source script
is not among the files reviewed when the no-adoption counterfactual was changed
to a common-steady-state comparison; the peak-$D^{G}$ anchor at $\varrho=0.8$
(9.18) does not match a fresh run of the current baseline calibration (8.87 pp,
verified). Regenerate this table --- ideally by writing a runner that emits it
via \texttt{latex\_tables.py} so it is auto-generated like the paper's other
exhibits --- before submission.}\par\smallskip
\begin{tabular}{lrrrr}
\toprule
$\varrho$ & $\psi_g$ (re-solved) & no policy peak $D^{G}$ & cap peak $D^{G}$ & transfer peak $D^{G}$\\
\midrule
0.60 & 0.714 & 8.39 & 0.03 & 8.53\\
0.70 & 0.624 & 8.81 & 0.02 & 8.95\\
0.80 & 0.538 & 9.18 & 0.01 & 9.32\\
\bottomrule
\end{tabular}
\caption{Robustness to the green--brown cost ratio $\varrho$ (price shock; peak
green-share response, pp over 24 quarters; $\psi_g$ re-solved each row to
$D^{G}_{ss}=0.05$). Cumulative output moves monotonically but the ordering does
not: no policy $-27.9\!\to\!-30.1$, cap $-21.1\!\to\!-20.9$, transfer
$+14.9\!\to\!+12.7$ as $\varrho$ rises from $0.6$ to $0.8$.}
\label{tab:rho}
\end{table}

\section{Experiments and results}

\subsection{Constructing the no-adoption counterfactual}
\label{sec:counterfactual}

Several results below isolate the adoption channel's contribution as the
difference between the full model and a "no-adoption" counterfactual,
$\Delta_H\equiv\sum_{t<H}(y^{\mathrm{adopt}}_t-y^{\mathrm{no\text{-}adopt}}_t)$.
The counterfactual is constructed to share the full model's steady state
\emph{exactly}: households still hold the discrete durable-choice margin, and
the steady-state green share, switching flow, and switching cost
($\bar D^{G}=5\%$, $D^{sw}_{ss}$, $\psi_g$) are bit-identical to the adoption
economy. What differs is the \emph{transition} only: the discrete-choice
probabilities are held at their steady-state values throughout the impulse
response (the choice stage's law of motion is the steady-state one, undifferentiated
by the shock), so the durable composition does not respond to the crisis, while
every other block --- prices, wages, the household's continuous
consumption/saving choice, monetary and fiscal policy --- responds normally. We
refer to this as \textbf{common-steady-state}, or \emph{frozen-choice}.

This matters because the alternative construction --- shutting the discrete
choice altogether by making green durables infeasible, so the counterfactual's
own steady state has $D^{G}_{ss}=0$ --- compares two economies that differ both
in the crisis response \emph{and} in steady-state composition and asset
holdings. We verified directly that this confound is quantitatively large, not
a second-order concern: at the baseline calibration, the mixed-steady-state
version of $\Delta_H$ for cumulative output under the price shock is smaller
than the common-steady-state version by roughly a factor of seven, because the
mixed comparison nets the crisis-time switching cost against an unrelated
steady-state cushioning effect of already having $5\%$ of households green at
date zero. All adoption-channel numbers in this paper use the common-steady-state
construction.


The \emph{price} shock is a 100 log-point rise in the world energy price with
a 16-quarter half-life, as in \citet{ARS2024}. The \emph{supply} shock is a 10\%
cut in available energy for 6 quarters, after which the world price clears
the market endogenously. Impulse responses are reported as $100\times$ the
level deviation from steady state. Variables whose steady state is one (prices,
$w$, $n$, and the population shares $D^{G}$, $D^{sw}$) read directly as percent
or percentage-point deviations, and $C$ (steady state $\approx0.997$) is within
rounding of a percent; $y$ and $\mathrm{gdp}$ have steady state
$\approx0.985$, so their entries are level deviations $\times100$, within about
1.5\% of the exact percentage. Energy quantities are a share of the energy
steady state ($\approx0.040$).

\input{output/tab_summary_core_import.tex}

\subsection{The crisis moves the transition, in a shock-dependent direction}

With no policy, the green share rises by up to \NEEDSRERUN{9.2} percentage points under
the price shock (from 5.0\% to 14.2\% of households) and \NEEDSRERUN{2.3} pp under the
supply shock (Figure~\ref{fig:baseline}). The output effect of the adoption
margin, however, is not one-signed. Under the \textbf{supply} shock it is
expansionary ($+\NEEDSRERUN{0.38}$ points of cumulative output with no policy, $+\NEEDSRERUN{1.67}$
under the transfer): adopters escape the energy bill and their real income
falls less. Under the \textbf{price} shock it is contractionary ($\NEEDSRERUN{-1.60}$
points with no policy, $\NEEDSRERUN{-4.65}$ under the transfer): the switching cost is
booked as an import in exactly the periods household income is falling
hardest (Section~\ref{sec:bop}), and at the calibrated $\psi_g$ this resource
cost outweighs the energy-bill saving. The sign of the adoption channel's
contribution to output is therefore not a general property of the model---it
depends on the calibrated switching cost, on which shock is driving
adoption, and (Section~\ref{sec:limits}) on the accounting of $\psi_g$.
\textbf{All $\Delta_H$ figures in this subsection predate the switch to the
common-steady-state counterfactual (Section~\ref{sec:counterfactual}) and must
be regenerated}; the signs above are independently verified to survive the
switch, but the magnitudes are understated by roughly a factor of seven at the
baseline calibration (see Section~\ref{sec:counterfactual}), so every number
above should be treated as illustrative pending rerun of \texttt{main.py}.

\begin{figure}[H]\centering
\includegraphics[width=0.92\textwidth]{output/fig1_price_shock.png}
\caption{Brown energy \emph{price} shock (elastic supply), adoption open
(solid) vs.\ shut (dashed), no policy.}
\label{fig:baseline}
\end{figure}

\begin{figure}[H]\centering
\includegraphics[width=0.92\textwidth]{output/fig1_supply_shock.png}
\caption{Brown energy \emph{supply} shock ($-10\%$ for 6 quarters, fixed
quantity), adoption open (solid) vs.\ shut (dashed), no policy.}
\label{fig:baseline_supply}
\end{figure}

\subsection{A sign condition for the adoption channel}
\label{sec:signmap}

The shock-dependence of the adoption channel's output contribution is not
accidental; it follows from a race between two resource flows, which we can
characterise and then verify. Write the contribution to cumulative output over
horizon $H$ as $\Delta_H\equiv\sum_{t<H}\big(y^{\text{adopt}}_t-y^{\text{no-adopt}}_t\big)$.
Adoption moves two things in the resource constraint (Section~\ref{sec:bop}):

\begin{itemize}\itemsep2pt
\item a \textbf{switching cost} $\psi_g D^{sw}_t$, booked as an import---a
resource \emph{outflow}, hence contractionary---which is \emph{front-loaded}
(paid once, up front, by each switcher) and scales with the adoption
\emph{volume} $D^{sw}$;
\item a \textbf{flow energy saving} $(P^{E}_{B}-P^{E}_{G})\,C^{G}_{E,t}$, booked
as reduced imports---a resource \emph{inflow}, hence expansionary---which
accrues \emph{gradually} over the life of the green durable.
\end{itemize}

\noindent The sign of $\Delta_H$ is the sign of the second flow, cumulated over
the window, minus the first. Two comparative statics follow.

\paragraph{Persistence (price closure).} A more persistent brown-price shock
raises the return to switching and therefore the adoption volume $D^{sw}$
steeply (peak $D^{G}$ rises from $1$ to $22$ pp as the half-life goes from $2$
to $256$ quarters). The front-loaded import cost thus grows faster than the
windowed saving, so $\Delta_H$ is \emph{decreasing} in persistence and crosses
zero at a finite half-life---\NEEDSRERUN{about $11$ quarters} in our calibration
(Figure~\ref{fig:signmap}). \textbf{This crossover point is the single number in
the paper most likely to move substantially under the common-steady-state
counterfactual} (Section~\ref{sec:counterfactual}): the common-SS $\Delta_H$ is
both larger in magnitude and decays far more slowly than the mixed-SS version
computed here (verified on the price shock: still an order of magnitude larger
than the mixed-SS channel at a 30-year horizon), which likely shortens the
half-life at which the front-loaded cost dominates. \emph{Do not cite the
11-quarter figure, or treat the sign of Route A's headline claim in
Section~\ref{sec:routeA} as verified, before rerunning
\texttt{run\_persistence.py} under the current code and re-checking this
number.} Short-lived shocks trigger little adoption whose
quick saving dominates ($\Delta_H>0$); persistent shocks trigger a large
adoption wave whose front-loaded cost dominates within the window
($\Delta_H<0$). This is the opposite of the naive intuition that more
persistence, by raising the per-switcher saving, should make adoption more
expansionary: the per-switcher calculus is dominated by the response of the
switcher \emph{count}.

\paragraph{Closure.} Under the fixed-quantity (supply) closure a third,
expansionary term appears that is absent under the price closure: adoption
lowers fossil demand against a fixed supply, so the market-clearing price
falls, relaxing the scarcity for \emph{all} households. This shifts $\Delta_H$
upward, which is why the supply shock's adoption contribution is positive
($+\NEEDSRERUN{0.38}$) at a persistence where the price shock's is negative.

\paragraph{Booking.} Both statics are conditional on the import booking of the
adoption flows. We solve the model under the alternative domestic booking of
Section~\ref{sec:bop}---green energy and green installation as competitive,
zero-profit domestic value added rebated to households, with only brown energy
imported---which resolves the phantom-dividend problem and clears to machine
precision. Under this booking neither the switching cost nor the green energy
bill is a leakage, and the adoption channel's contribution to cumulative output
($\Delta_{24}$) flips sign:
\IfFileExists{output/tab_signmap_booking.tex}{\begin{table}[H]\centering
\small
\begin{tabular}{lcc}
\toprule
$\Delta_{24}$, adoption output contribution & import booking & domestic booking\\
\midrule
price shock & $-1.60$ & $+35.98$\\
supply shock & $+0.38$ & $+5.96$\\
\bottomrule
\end{tabular}
\caption{Adoption channel contribution to cumulative output ($\Delta_{24}$), import vs.\ domestic booking (Section~\ref{sec:bop}). \emph{None} policy.}
\label{tab:signmap}
\end{table}
}{\textit{[table \texttt{tab\_signmap\_booking}: run \texttt{run\_signmap.py}]}}
\textbf{This table is auto-generated and will pick up the common-steady-state
counterfactual automatically on rerun} (Section~\ref{sec:counterfactual}) ---
\emph{provided} \texttt{run\_signmap.py} calls the shared \texttt{run(...,
model\_variant=)} API rather than reimplementing the steady-state comparison
itself; verify this before trusting the regenerated numbers. The domestic-booking
combination has not been separately checked against the common-SS construction.
The adoption wave itself is essentially unchanged across bookings (peak $D^{G}$
$\NEEDSRERUN{9.18\to9.46}$ pp on the price shock, $\NEEDSRERUN{2.33\to1.34}$ on the supply shock): only
the booking of its resource flows moves, not the margin. The domestic-booking
magnitudes are an \emph{upper bound}---green value added is modelled as costless,
maximally stimulative---but the \emph{sign reversal} is robust and confirms the
characterisation: the sign of the adoption channel is a property of how its
resource cost is booked, not of the adoption margin itself. Finally, the marginal switcher's steady-state indifference,
$\psi_g\approx(\bar P^{E}_{B}-\bar P^{E}_{G})\,\bar C^{G}_{E}/(r+\delta_g)$,
ties the threshold to the cost ratio $\varrho$ and reproduces the direction of
the $\psi_g(\varrho)$ mapping in Table~\ref{tab:rho}. A closed-form threshold
via the full Keynesian-cross algebra of the adoption block is left for the next
draft; the characterisation above is verified numerically in
Figure~\ref{fig:signmap}.

\begin{figure}[H]\centering
\IfFileExists{output/fig_signmap.png}{\includegraphics[width=0.68\textwidth]{output/fig_signmap.png}}{\textit{[figure output/fig\_signmap.png: run run\_signmap.py]}}
\caption{Adoption's contribution to cumulative output, $\Delta_{24}$, against
the brown-price-shock half-life (price closure, orange), with the supply-shock
value marked (blue). The contribution is expansionary for transitory shocks,
crosses zero near an $11$-quarter half-life, and turns increasingly
contractionary as the shock---and hence the adoption wave whose front-loaded
switching cost it imports---becomes more persistent.}
\label{fig:signmap}
\end{figure}

\subsection{The price cap eliminates the adoption margin}
\label{sec:cap}

Under the cap, the peak green share response falls to essentially zero under
both shocks (from $\NEEDSRERUN{+9.18}$ to $+0.01$ pp on the price shock, from $\NEEDSRERUN{+2.33}$ to
$+0.01$ on the supply shock): the adoption channel is not attenuated, it is
switched off. With the margin shut its contribution to cumulative output
collapses toward zero accordingly (this last point is the $\Delta_H$ comparison
and should be re-verified under the common-steady-state counterfactual,
Section~\ref{sec:counterfactual}, though the direction is not in doubt: the cap
holds $D^{G}$ at its steady-state value, which is exactly the frozen
counterfactual's construction, so under the cap the two are close by
construction). The mechanism is direct: the cap holds
$P^{E}_{B}$ at its pre-crisis level, and $P^{E}_{B}$ is precisely what determines
the return to adoption.

\begin{figure}[H]\centering
\includegraphics[width=0.92\textwidth]{output/fig2_policy_price.png}
\caption{Fiscal response to the \emph{price} shock, adoption open: no policy
(solid), price cap (dashed), Slutsky transfer (dotted).}
\label{fig:policy_price}
\end{figure}

\begin{figure}[H]\centering
\includegraphics[width=0.92\textwidth]{output/fig2_policy_supply.png}
\caption{Fiscal response to the \emph{supply} shock, adoption open: no policy
(solid), price cap (dashed), Slutsky transfer (dotted).}
\label{fig:policy_supply}
\end{figure}

\subsection{The transfer preserves adoption, at a cost to output}
\label{sec:transfer}

The Slutsky transfer leaves the green share response essentially unchanged
($\NEEDSRERUN{9.46}$ vs.\ $\NEEDSRERUN{9.18}$ pp under the price shock, $\NEEDSRERUN{3.35}$ vs.\ $\NEEDSRERUN{2.33}$ under the
supply shock) while delivering more output stabilisation than the cap, at
similar fiscal cost: subsidies destroy the transition margin, transfers do
not. This preservation is a direct consequence of the transfer's lump-sum
indexing (Section~\ref{sec:bop}): because it does not condition on the
household's current durable, switching forfeits nothing and the relative-price
incentive to adopt is left intact. Leaving adoption intact is not free for output under the price shock,
however: the transfer's cumulative-output gain is \emph{smaller} with
adoption open than under the common-steady-state counterfactual
($\NEEDSRERUN{-4.65}$ points from the margin, Table~\ref{tab:summary}, which is
auto-generated and will update on rerun),
for the same resource-cost reason as above.

\subsection{Fixed supply makes the cap destructive}

Under the supply shock the cap yields cumulative output of $-46.9$ against
$-8.2$ with no policy: with inelastic supply, capping the price prevents
substitution away from energy and pushes the economy towards Leontief. This
survives the addition of an adoption margin. It is also precisely the
configuration in which \citet{Langot2026}, assuming elastic supply, find the
cap effective: our elastic-supply column agrees with them ($-20.9$ vs.\
$-30.1$ for cap vs.\ no policy under the price shock).

\begin{figure}[H]\centering
\includegraphics[width=\textwidth]{output/fig3_policy_adoption.png}
\caption{Policy and the adoption margin. Under both shocks the price cap
(dashed) flattens the green share response to zero; the transfer (dotted)
leaves it intact.}
\end{figure}

\subsection{The cap is a dial, not a switch}

Table~\ref{tab:summary} used $\tau^{E}=1$, a complete price freeze. Table~\ref{tab:dose}
sweeps $\tau^{E}\in[0,1]$. Under the \emph{price} shock the adoption loss is
almost exactly linear in $\tau^{E}$ (half a cap costs a bit under half the
transition: $\NEEDSRERUN{4.53}$ against $\NEEDSRERUN{9.18}$ pp). Under the \emph{supply} shock it is convex:
small caps barely touch adoption and the collapse is concentrated near the
full freeze. A partial shield is therefore much less damaging to the
transition when supply is inelastic---but, as the next result shows, much
more damaging to everything else.

\input{output/tab_dose_core_import.tex}

\subsection{The welfare effect of the cap flips sign with the supply elasticity}

This is the sharpest result in the paper. Under the price shock, tightening
the cap \emph{raises} welfare monotonically, from $-1.30\%$ to $-0.57\%$;
under the supply shock it \emph{lowers} welfare monotonically, from $-0.59\%$
to $-1.28\%$. The sign of the cap's welfare effect therefore hinges on the
supply elasticity. \citet{Langot2026}, assuming elastic supply, find the French
shield effective; \citet{Bayer2026}, assuming fixed union-wide supply, find
subsidies destructive. Our two closures reproduce the \emph{sign} each obtains,
which locates one source of the contrast in a single parameter. We do not claim
to reconcile the two papers fully: Langot's result also runs through a
labour-cost channel and Bayer's through cross-border spillovers in a monetary
union, and our small open economy nests neither. The narrower point we
establish is that, once the supply elasticity is fixed, the direction of the
cap's welfare effect is determined---so part of the empirical disagreement
reduces to measuring that elasticity at the relevant horizon.

\begin{figure}[H]\centering
\includegraphics[width=\textwidth]{output/fig6_dose_response.png}
\caption{Dose-response in $\tau^{E}$ (orange) against the Slutsky transfer
benchmark (dotted). Note the sign reversal of the welfare panel between the
two shocks.}
\end{figure}

\paragraph{Private welfare and the transition are distinct.} The CEV above is a
\emph{private} household-welfare measure: it values consumption and leisure
along the transition and prices no climate or transition externality. Two
results should therefore be kept separate. First, on private welfare alone the
Slutsky transfer dominates the cap under \emph{both} shocks ($-0.25$ vs.\
$-0.57$ under the price shock, $-0.16$ vs.\ $-1.28$ under the supply shock,
Table~\ref{tab:dose}); the case for the transfer over the cap does not rest on
the transition. Second, and separately, the transfer preserves the green
adoption margin that the cap eliminates (Section~4.3)---a positive statement
about the energy transition, not a welfare ranking. We deliberately do not
aggregate the two into a single social-welfare number: that would require a
social value of adoption (a carbon price times avoided brown-energy use) that
we do not model. A positive such value only reinforces the private ranking; it
cannot reverse it here, since the transfer already dominates.

\subsection{Distributional incidence}

\input{output/tab_cev_core_import.tex}

Without policy the impatient lose far more than the patient ($-2.10$ vs.\
$-0.14$, roughly $15\times$). The patient remain the least-hurt group
throughout the cap sweep. One clear distributional result: under the
transfer, the \textbf{patient type's CEV turns positive} ($+0.17\%$) while
impatient and middle types remain negative. Indexed to pre-crisis energy use,
the Slutsky compensation over-compensates low-MPC, low-energy-exposure
households relative to their true loss---a symptom of the fiscal block's
lack of investment margin to crowd out the transfer (Section~\ref{sec:limits}).

% =====================================================================
% Section 5 -- Ex ante greening versus ex post insurance
% Paste into ehank_results.tex after the Experiments section (Sec. 4)
% and before Relation to the literature. All numbers below are the
% taste_shock = 0.05 / import-booking calibration; the \input-ed tables
% and figures regenerate them from run_summary_table.py,
% run_exante_expost.py and run_persistence.py.
% =====================================================================

\subsection{The brown-anchored CPI understates the cost of living}
\label{sec:deflator}

The model's CPI anchors its energy leg on the brown price alone
(eq.~\ref{eq:innernest}), while the economically correct index reweights brown
and green by the adoption share. The gap has the closed form
$\phi^{1-\eta_E}=1+\alpha_E D^{G}\bigl[(P^E_G)^{1-\eta_E}-(P^E_B)^{1-\eta_E}\bigr]$,
computed ex post so it closes no loop in the model. It is small but not zero:
the measured deflator overstates the cost of living by at most a few tenths of
an annualised percentage point over the crisis, and the overstatement grows
with adoption (Figure~\ref{fig:deflator}). Because the monetary rule reads the
measured CPI, this is the wedge a Taylor rule would miss; under the
domestic-good numeraire it affects only the measured deflator and the rule, not
the resource constraint.

\begin{figure}[H]\centering
\includegraphics[width=0.85\textwidth]{output/fig8_deflator_gap.png}
\caption{True minus measured annualised inflation from anchoring the CPI on
$P^E_B$ alone, by policy.}
\label{fig:deflator}
\end{figure}


\subsection{Robustness: the nonlinearity lives in the adoption margin}
\label{sec:nonlin}

A natural concern is that the discrete adoption choice makes the economy too
nonlinear for the linearised (sequence-space) solution to be reliable, in
contrast to \citet{ARS2024}, who find nonlinearity immaterial in the continuous
economy. Comparing the linearised impulse responses with their nonlinear
counterparts (a fixed--Jacobian Newton solve) shows the concern is confined to
one variable. Aggregate output is essentially linear: the ratio of the
nonlinear to the linear peak response, $y^{\mathrm{NL}}/y^{\mathrm{L}}$, stays
in $[1.00,1.01]$ across shock sizes, reproducing the small nonlinearity ARS
report. The nonlinearity is concentrated in the adoption share $D^{G}$, whose
ratio rises from $\NEEDSRERUN{1.16}$ at a one-eighth shock to $\NEEDSRERUN{2.04}$ at a three-quarter
shock (Figure~\ref{fig:nonlin}, left). With the margin \NEEDSRERUN{shut}, the bare household
block is in fact mildly \emph{concave} ($y^{\mathrm{NL}}/y^{\mathrm{L}}\approx
\NEEDSRERUN{0.92}$ at the largest shock); the adoption margin is \emph{convex} and partly
offsets it, which is why output is, if anything, more linear with the margin
open. \textbf{The "margin shut" comparison in this subsection predates the
common-steady-state counterfactual} (Section~\ref{sec:counterfactual}) and uses
the OLD, zero-steady-state construction, which is arguably still the right
comparison for THIS specific exercise (isolating the bare continuous-HANK
block's own curvature, to compare against ARS's technology-free economy, is a
different question from isolating the crisis response of the discrete choice);
but for consistency with the rest of the paper, and because \texttt{model.py}
no longer exposes that construction under \texttt{model\_variant='no\_adoption'}, this comparison
should be redone against the frozen (common-SS) counterfactual before
submission, with an explicit note on which of the two questions is being asked.

The convexity is the intrinsic curvature of a logit share evaluated near a low
level. For a binary choice $P=[1+\exp(-\Delta/\sigma_\varepsilon)]^{-1}$ in the
value gap $\Delta$, the relative curvature is $(1-2P)/\sigma_\varepsilon$, so
$D^{G,\mathrm{NL}}/D^{G,\mathrm{L}}-1\propto 1/\sigma_\varepsilon$. Re-solving on
a grid of $\sigma_\varepsilon$ (holding $\bar D^{G}=0.05$ by recalibrating
$\psi_g$) confirms this: the ratio falls monotonically from $2.21$ at
$\sigma_\varepsilon=0.02$ to $1.06$ at $\sigma_\varepsilon=0.40$
(Figure~\ref{fig:nonlin}, right), and it tends to one as the shock vanishes, so
the linearisation is exact to first order and the Jacobian is correct. Two
implications follow. Aggregate and consumption-based welfare statistics are
reliable under linearisation. The \emph{level} of induced greening at the
largest shocks is understated by the linear solution---and does not converge
under a full-size shock---so adoption-channel magnitudes are read at a reduced
shock size, and the taste scale $\sigma_\varepsilon$ that governs this curvature
is the same parameter left underidentified in Section~\ref{sec:limits}: an
adoption-elasticity moment would pin both at once.

\begin{figure}[H]\centering
\includegraphics[width=\textwidth]{output/fig9_nonlinearity.png}
\caption{The nonlinearity is confined to the discrete adoption margin.
\emph{Left:} peak nonlinear/linear ratio vs shock size---output (adoption on
and off) and $D^{G}$. \emph{Right:} the same ratio vs the logit taste scale
$\sigma_\varepsilon$ at a half-size shock, $\psi_g$ recalibrated to
$\bar D^{G}=5\%$.}
\label{fig:nonlin}
\end{figure}

\section{Ex ante greening versus ex post insurance}
\label{sec:exante}

The instruments in Section~\ref{sec:cap}--\ref{sec:transfer} are all
\emph{ex post}: they are deployed once the price shock has landed. Because the
adoption margin responds to the same relative price the crisis moves, a natural
alternative is to act \emph{ex ante}---to lean on the margin in the steady state
with a permanent carbon tax $\tau_b$ on brown energy, leaving the economy greener
and, one might hope, less exposed when the shock arrives. This section prices
that hope. The answer is that ex ante greening is a poor substitute for ex post
insurance, and the reason is intrinsic to the channel this paper adds: the
adoption margin is an endogenous shock absorber, and pre-greening spends it.

Table~\ref{tab:summary_price_import} collects the instruments on one axis. The
welfare metric is the total consumption-equivalent variation (CEV) of each full
scenario relative to the common no-tax steady state, by discount-factor type and
averaged, computed ex post from the cached impulse response
(\texttt{welfare.cev\_total}); a permanent policy enters through the standing gap
between its steady state and the baseline, a crisis-time policy through the
transition alone, so the two are directly comparable.

\input{output/tab_summary_import.tex}

\subsection{Targeting energy transfers is regressive}
\label{sec:flat}

The Slutsky transfer of Section~\ref{sec:transfer} indexes the lump sum to each
household's pre-crisis brown-energy use. Replacing it with an
\emph{untargeted} transfer of the \emph{same aggregate envelope}---the identical
outlay ($49.3$ vs.\ $49.3$ over $H=24$), paid flat across the distribution rather
than indexed---almost halves the welfare loss, from a CEV of $-0.82\%$ to
$-0.32\%$ (Table~\ref{tab:summary_price_import}). The gain is entirely
distributional. By type, the flat transfer moves the impatient (high-MPC)
household from $-0.68\%$ to $+0.10\%$---it is made \emph{better off} than in the
no-shock steady state---while the patient household bears more of the burden
($-0.77\%$); the Slutsky transfer, by contrast, spreads the loss almost evenly
($-0.68$, $-0.89$, $-0.88$).

The mechanism is that energy-indexed targeting is regressive \emph{in this
model}. With homothetic CES energy demand the energy share is constant across the
wealth distribution, so brown-energy use rises with total consumption, hence with
permanent income and patience. Indexing the transfer to pre-crisis energy use
therefore hands the largest cheques to the low-MPC, wealthy households, exactly
where a euro buys the least insurance. A flat transfer reaches the high-MPC
households the crisis hits hardest. This speaks directly to the design of the
2022--23 European energy-price ``brakes,'' which were indexed to historical
consumption: within the model, the indexing is the source of the inefficiency,
not a refinement of it.

We flag the caveat plainly. The result rests on the homothetic energy nest: if
energy were a necessity with a declining Engel curve---the empirically standard
case---low-income households would spend a \emph{larger} share on energy and the
sign could reverse. The clean statement is conditional: \emph{under constant
energy shares, energy-indexed targeting is dominated by an equal-value flat
transfer}. Making the nest non-homothetic is the right robustness check and is
left to the calibration section's agenda.

\subsection{A welfare-optimal permanent carbon tax}
\label{sec:carbontax}

Before comparing ex ante and ex post, we characterise the steady state a
permanent $\tau_b$ delivers. Holding $\psi_g$ fixed at its no-tax value and
letting $D^{G}_{ss}$ float, a carbon tax raises the effective brown/green price
ratio and lifts steady-state adoption: $D^{G}_{ss}$ moves from $5\%$ at
$\tau_b=0$ to $8.9\%$ at $\tau_b=0.10$ and $20.1\%$ at $\tau_b=0.30$
(Figure~\ref{fig:carbon_optimal}, right). The carbon revenue is rebated
lump-sum, so the exercise is budget-neutral and isolates the allocative effect.

Steady-state welfare is hump-shaped in $\tau_b$
(Figure~\ref{fig:carbon_optimal}, left): the standing CEV rises to a peak of
$+0.043\%$ at $\tau_b^{*}\approx 10\%$ and turns negative beyond
$\tau_b\approx 25\%$. The model contains \emph{no} climate externality, so this
interior optimum is not Pigouvian. It is a second-best correction of the
adoption friction: green energy is structurally cheaper
($\bar P^{E}_{G}=0.8<\bar P^{E}_{B}=1$) but the logit taste-shock scale
$\sigma_\varepsilon$ and the switching cost $\psi_g$ leave green under-adopted, so
a modest tax that pushes households toward the cheaper technology raises welfare
until the tax distortion dominates. That the optimum is positive at all is a
property of this friction, and we treat it as a diagnostic of the channel rather
than a policy recommendation---it would be swamped by any explicitly modelled
externality or abatement cost.

\begin{figure}[H]\centering
\includegraphics[width=\textwidth]{output/fig_carbon_optimal_import.pdf}
\caption{The welfare-optimal permanent carbon tax. Left: standing CEV of the
$\tau_b$ steady state relative to the no-tax baseline (no crisis), hump-shaped
with an interior optimum near $\tau_b^{*}\approx10\%$. Right: the induced
steady-state green share $D^{G}_{ss}$. $\psi_g$ is held fixed and $D^{G}_{ss}$
floats; the carbon revenue is rebated lump-sum.}
\label{fig:carbon_optimal}
\end{figure}

\subsection{Ex ante greening barely insures against the crisis}
\label{sec:versionA}

We now hit two economies with the \emph{same} price shock: the baseline
($D^{G}_{ss}=5\%$) and a ``prepared'' economy carrying the welfare-optimal
$\tau_b=0.10$ ($D^{G}_{ss}=8.9\%$). Table~\ref{tab:exante_expost_import}
reports the total CEV of each, alongside the ex post cap and transfer on the
baseline economy.

\input{output/tab_exante_expost_import.tex}

The prepared economy is essentially no better insured than the unprepared one:
its crisis CEV is $-1.85\%$ against $-1.87\%$ for laissez-faire, a difference of
two basis points, and both are far worse than either ex post instrument (cap
$-1.15\%$, transfer $-0.82\%$). Adding the standing benefit of the tax
($+0.04\%$, Section~\ref{sec:carbontax}) does not change the ranking. The
greening the tax buys is real---$D^{G}_{ss}$ is $78\%$ higher---but too small a
share of the economy to move aggregate exposure appreciably, and the tax's
crisis-time value is a rounding error next to a transfer that puts cash in
high-MPC hands exactly when it binds. Consistent with Section~\ref{sec:cap}, the
combination ``ETS $+$ cap'' tracks the cap alone: once the price is frozen the
adoption margin is shut and the greener starting point is irrelevant.

The green \emph{subsidy} of Section~\ref{sec:cap}\footnote{The transitory
switching subsidy $s_g$, financed by debt over the crisis window.} is worth a
line as a crisis instrument in its own right. It doubles the adoption response
(peak $D^{G}$ of $18.0$ vs.\ $8.9$ pp) at trivial fiscal cost, yet its CEV is
the worst in the table, $-2.03\%$, below laissez-faire. Paying households to
incur the real switching cost $\psi_g D^{sw}$---booked as an import
(Section~\ref{sec:bop})---in the middle of a consumption crisis destroys welfare
even as it greens the economy. The adoption margin is valuable, but not as
something to force at the trough.

\subsection{Persistence exhausts the adoption absorber}
\label{sec:routeA}

The obvious defence of the ex ante case is persistence: a longer-lived shock
should let the greener economy's lower exposure compound over more periods.
Figure~\ref{fig:persistence} and Table~\ref{tab:persistence_import} test it by
sweeping the shock half-life and re-solving only the impulse response (the steady
state does not move with persistence), decomposing the ex ante number into its
constant standing part and a \emph{crisis-only greening dividend}, defined as the
prepared economy's crisis CEV minus the standing $+0.04\%$ minus laissez-faire.

\IfFileExists{output/tab_persistence_import.tex}{\input{output/tab_persistence_import.tex}}{\textit{[table \texttt{tab\_persistence\_import}: run \texttt{run\_persistence.py}]}}

The defence fails, and in an instructive way. The greening dividend does not grow
with persistence---it \emph{turns negative}, from $+0.01\%$ at a four-quarter
half-life to $-0.09\%$ at a twelve-year one, crossing zero near a three-year
half-life. The ex post cap moves the other way: its crisis protection rises
monotonically, from $+0.24\%$ to $+1.36\%$, as a longer freeze shields more of
the transition.

The reason is the adoption margin itself. A persistent brown-price shock makes
switching attractive for many quarters, so the \emph{browner} laissez-faire
economy adopts heavily \emph{during} the crisis and reaps that windfall over the
whole episode---the margin is an endogenous shock absorber. The prepared economy
has already spent much of that absorber in the steady state: with fewer brown
households left at the switching margin, it has less adoption to do when the shock
arrives, and less of the crisis-time windfall to collect. Pre-greening does not
merely fail to insure against a persistent shock; it \emph{consumes the very
margin} that would have cushioned it. This is the sharpest form of the paper's
thesis. The instruments that freeze the adoption technology---the cap here, and
by construction the frozen-technology models of \citet{Bayer2026},
\citet{Langot2026} and \citet{ARS2024}---miss an absorber that works best when
left loaded, i.e.\ when the economy enters the crisis with room to green rather
than already green.

\begin{figure}[H]\centering
\IfFileExists{output/fig_persistence_import.pdf}{\includegraphics[width=\textwidth]{output/fig_persistence_import.pdf}}{\textit{[figure output/fig\_persistence\_import.pdf: run run\_persistence.py]}}
\caption{Persistence and the two policies. Left: total CEV of laissez-faire, the
ex post cap and the ex ante ETS as the shock half-life rises. Right: the
crisis-only gain over laissez-faire---the cap's protection grows with
persistence, while the ex ante \emph{greening dividend} turns negative, because a
more persistent shock lets the browner economy exploit the adoption absorber the
prepared economy has already spent.}
\label{fig:persistence}
\end{figure}

\paragraph{Takeaway.} Three results point the same way. Energy-indexed targeting
is dominated by a flat transfer of equal budget (Section~\ref{sec:flat}); a
permanent carbon tax is mildly welfare-improving in the steady state through a
non-Pigouvian adoption correction, but essentially worthless as crisis insurance
(Sections~\ref{sec:carbontax}--\ref{sec:versionA}); and its insurance value turns
negative as shocks become persistent, because greening ex ante exhausts the
adoption absorber (Section~\ref{sec:routeA}). The margin this paper adds is not a
lever a planner should pre-pull; it is a stabiliser that pays off precisely when
the economy is left free to adopt in the crisis.



\subsection{Recycling the carbon revenue: rebate versus green subsidy}
\label{sec:recycle}

The welfare-optimal carbon tax of Section~\ref{sec:carbontax} returns its
revenue lump sum. An alternative earmarks it to the transition itself, paying a
share $s_g$ of the green switching cost. We compare the two recyclings across
carbon-tax rates $\tau_b\in\{0.10,0.20,0.30\}$, holding $\psi_g$ at its no-tax
value so the steady-state green share floats (Figures~\ref{fig:taub_rebate}
and~\ref{fig:taub_green}). Under the lump-sum rebate the steady-state green
share rises from $8.9\%$ to $20.1\%$ across the grid, and the whole brown-price
path in the crisis shifts up by the factor $(1+\tau_b)$: the pre-tax price
passes the world shock through identically, while households pay the
tax-inclusive price. Under \emph{budget-neutral} green-subsidy recycling---$s_g$
solved so that the switching subsidy exactly exhausts the carbon revenue
($T^{\mathrm{reb}}=0$)---the revenue funds a subsidy covering $35$--$40\%$ of
the switching cost and lifts the steady-state green share to $27$--$46\%$.

Two features are worth flagging. First, the carbon base is self-eroding: at
$\tau_b=0.10$ the revenue is a third smaller under green-subsidy recycling than
under the rebate, because the greener economy holds less brown-durable stock to
tax. Second, the crisis-time greening response is \emph{smaller}, not larger,
from the greener starting point under the subsidy: recycling that greens the
economy through a lower switching \emph{cost} preferentially converts the
households nearest the logit margin, leaving a residual brown pool that is less
responsive to the shock---the mirror image of the rebate, which greens through a
higher permanent brown \emph{price} and leaves the residual pool closer to the
margin. This is the adoption-absorber exhaustion of Section~\ref{sec:routeA}
seen across policy instruments rather than across shock persistence.

\begin{figure}[H]\centering
\includegraphics[width=\textwidth]{output/fig_exante_taub_rebate.png}
\caption{Carbon-tax sweep, lump-sum rebate: laissez-faire and ex-post cap
against the ex-ante ETS economy at $\tau_b\in\{0.10,0.20,0.30\}$.}
\label{fig:taub_rebate}
\end{figure}

\begin{figure}[H]\centering
\includegraphics[width=\textwidth]{output/fig_exante_taub_greensub_bn.png}
\caption{Carbon-tax sweep, budget-neutral green-subsidy recycling
($T^{\mathrm{reb}}=0$, $s_g$ endogenous). Same axes as
Figure~\ref{fig:taub_rebate}.}
\label{fig:taub_green}
\end{figure}

\section{Relation to the literature}
\label{sec:lit}

This section is written for the coauthors: it states, paper by paper, what our
setup changes relative to the reference and what that change does to the
results. Compact one-page summaries of each reference are in
Appendix~\ref{sec:summaries}.

\subsection{Relative to \citet{ARS2024}}

\paragraph{Setup.} We inherit the entire ARS aggregate economy---the CES
energy nest, sticky import prices, the mutual fund, the energy suppliers with
intertemporal arbitrage, the real-wage-stabilising wage Phillips curve, and the
constant-real-rate monetary baseline---and change the household in four coupled
ways. (i) ARS's household is homothetic with a fixed energy technology: every
household shares one CES basket and one CPI, and energy demand is
$\alpha_E(P_E/P)^{-\eta_E}c$. We add a discrete brown/green durable state and a
DC-EGM adoption choice, so the energy technology becomes an endogenous,
household-level object. (ii) ARS has a single household energy price; we have
two ($P^E_B$, $P^E_G$) and a type-specific cost-of-living index. (iii) ARS uses
the CPI as unit of account; we use the domestic good, so that the resource
constraint is invariant to how the CPI weights brown and green energy. (iv) We
add two current-account terms absent in ARS---the switching cost as an import,
the adopters' energy saving as an import saving. The energy-supply closure is a
generalisation of ARS: their baseline is our $\varepsilon^E=\infty$
price-taking case, and their coordinated world equilibrium (endogenous world
price) is the two-country analogue of our $\varepsilon^E$-finite case.

\paragraph{Results.} ARS's central mechanism---low $\chi$ makes the energy
price shock contractionary through the real-income channel---is preserved and
operates in the background of every one of our experiments. Where we depart is
that ARS's fiscal ranking is entirely about output, inflation and inequality:
in ARS the subsidy is a domestic ``silver bullet'' that stabilises all three,
and its only drawback is a \emph{cross-country} externality (coordinated
subsidies spike the world price and become self-defeating). Our model adds an
orthogonal cost of the same subsidy that ARS cannot see: capping $P^E_B$
eliminates the green adoption margin (peak $D^G$ from \NEEDSRERUN{9.18} to 0.01 pp), so the
subsidy is not neutral for the transition. Two further connections are worth
flagging to the coauthors. First, our supply-shock result---the cap is
destructive under inelastic supply---is the \emph{single-country} counterpart
of ARS's coordination result: in both, capping the price under an
endogenous/scarce energy price prevents the necessary substitution and deepens
the recession. Second, our monetary baseline being ARS's constant-real-rate
rule is what pins the real exchange rate, so a domestic transfer has no
terms-of-trade effect; this is inherited, not assumed anew.

\subsection{Relative to \citet{Bayer2026}}

\paragraph{Setup.} Bayer et al.\ is a two-country HANK monetary union (Germany
one-third, rest of the union two-thirds), with a fixed nominal exchange rate,
built to study \emph{spillovers}. Our economy is a single small open economy
with an exogenous foreign block, so by construction we cannot represent their
central object---the cross-border, beggar-thy-neighbour transmission of a
national subsidy. Three further setup differences matter. (i) Supply: Bayer
fix the union-wide energy \emph{quantity} and let a common price clear; this is
exactly our $\varepsilon^E$-finite closure, so our supply-shock experiments are
the natural comparison. Our shock is smaller (a 10\% cut for 6 quarters against
their 20\% cut for 6 quarters), chosen for first-order reliability. (ii) Supply
side: Bayer has a rich medium-scale block---two assets (liquid bonds, illiquid
capital), an investment margin with capacity utilisation, energy in production
as well as consumption, and GHH preferences. We have a single asset, no
capital, and energy in consumption only in the baseline. (iii) Heterogeneity:
Bayer's headline distributional mechanism is \emph{within-country
energy-share heterogeneity}---poor households have systematically higher energy
shares (a high/low-intensity Markov type correlated with income), so they
suffer larger real-income losses. We assume homothetic energy demand
conditional on the durable type, so we cannot speak to their energy-exposure
distributional results; our heterogeneity is over discount factors and over the
adoption margin instead. Finally, their monetary policy is a standard Taylor
rule, not the constant-real-rate rule; they show this matters, because under a
Taylor rule the subsidy also suppresses \emph{intertemporal} substitution
through household-specific real rates.

\paragraph{Results.} Under the supply closure we reproduce the qualitative
core of Bayer et al.: the Slutsky/Hicks transfer welfare-dominates the price
cap (our private CEV $-0.16$ vs.\ $-1.28$), and capping the price under a fixed
quantity is inefficient because it prevents the substitution the scarcity
requires---our ``cap pushes towards Leontief'' is their ``subsidising under a
supply contraction is inefficient.'' What we add is the adoption margin, which
Bayer et al.\ do not have: the same cap that is inefficient for output also
switches off the transition, and the transfer that is welfare-superior also
preserves it. What we \emph{cannot} deliver is their spillover result: with no
second country, we have nothing to say about whether one country's subsidy
harms its neighbours. This is the reason we are careful (Section~4.6, and the
abstract) not to claim that a single parameter ``reconciles'' Bayer and Langot;
we reproduce the domestic sign, not the cross-border mechanism.

\subsection{Relative to \citet{Langot2026}}

\paragraph{Setup.} Langot et al.\ study France as a single economy inside the
euro area, with an ECB Taylor rule down-weighted for France's size. Two
structural choices distinguish their model from ours. (i) They take the energy
price as \emph{exogenous} (fed from the government's official forecasts), so
France is a price taker facing a perfectly elastic supply at that price---this
is our $\varepsilon^E=\infty$ closure, and their price-shock scenario is the
right comparison. (ii) Their energy substitution elasticity is $\eta_E=0.5$,
five times ours, and they add an \emph{incompressible} (subsistence) level of
energy consumption, which delivers non-homotheticity: poor households have a
larger incompressible share, hence a higher effective energy share and a lower
price elasticity. This is their distributional mechanism, analogous in spirit
to Bayer's energy-share heterogeneity but generated by subsistence rather than
by types. Our model buys energy substitution on a different margin: we keep the
\emph{intensive} elasticity low ($\eta_E=0.1$, as in ARS) and obtain
substitution from the \emph{extensive} adoption margin instead. A third,
methodological difference matters for how the results should be read: Langot et
al.\ conduct an \emph{ex-ante, conditional-forecast} evaluation---they back out
the sequence of structural shocks that makes the HANK reproduce the French
government's published forecasts for output, inflation and debt, then run
counterfactual policies holding those shocks fixed. Ours (like ARS and Bayer)
is a clean MIT-shock impulse-response exercise, not a fit to a specific
historical episode.

\paragraph{Results.} Under the elastic (price) closure we agree with Langot's
headline: the cap cushions output ($\sum y = -20.9$ vs.\ $-30.1$ with no
policy) and, on private welfare, tightening the cap raises the CEV
($-1.30\to-0.57$)---consistent with their finding that the shield supports
growth and curbs inflation. But two qualifications are important for the
coauthors. First, Langot's pro-shield result rests substantially on a
\emph{labour-cost} channel: by containing the retail energy price the shield
prevents the price--wage spiral, holds down the cost of labour, and thereby
supports employment; in their comparison the shield beats both wage indexation
and a targeted transfer precisely on this employment margin. Our wage block and
our lack of an investment/employment margin mean we reproduce the \emph{sign}
of their result but not the full mechanism---hence we do not claim the
Langot--Bayer contrast is ``only'' about the supply elasticity. Second, Langot
et al.\ hold the energy technology fixed, so like ARS and Bayer they do not see
the cap's cost to the transition, which is our contribution.

\section{Limitations}
\label{sec:limits}

\begin{enumerate}\itemsep3pt

\item \textbf{The taste scale $\sigma_{\varepsilon}$ is calibrated, not
estimated.} $\psi_g$ and $\sigma_\varepsilon$ jointly hit $D^G_{ss}=0.05$, but
one level moment cannot separately identify both a switching cost and a taste
scale: a second moment (e.g.\ the semi-elasticity of adoption to a purchase
subsidy) is needed and would discipline the sign found in
Section~4.1 along with the magnitudes throughout. The \emph{ordering} of the
cap-vs-transfer results does not depend on it, since it follows from the
relative price alone.

\item \textbf{The switching cost is booked as an import by accounting
necessity.} Under ARS's $\mu_{ss}>1$ with capitalised profits, booking
$\psi_g D^{sw}$ as domestic demand creates a marginal dividend owned by nobody
and breaks asset-market clearing; we therefore book it externally
(Section~\ref{sec:bop}). This convention is what makes the adoption channel
\emph{contractionary} under the price shock, since the cost lands in the
periods household income is falling hardest (Section~4.1). The \emph{ordering}
of the policy results does not depend on it---it follows from the relative
price alone---but the sign of the adoption channel's output contribution under
a price shock does, and should be read as conditional on this booking.

\item \textbf{Zero pass-through to the green price} ($\rho_g=0$) means
adopters are completely insulated from the fossil shock: an upper bound on
the adoption channel. A partial pass-through calibrated to observed
electricity--gas co-movement would be more defensible.

\item \textbf{No modelled green technology.} The adopters' resource saving
rests on the domestic/imported interpretation of Section~\ref{sec:bop} rather
than on a green sector with its own capacity, cost curve and investment. We
have implemented the minimal version---a competitive, zero-profit installer
sector that books the switching cost as domestic value added rebated to
households---and it clears to machine precision and flips the sign of the
adoption channel (Section~\ref{sec:signmap}); what remains unmodelled is the
green energy \emph{supply} technology (capacity, cost curve, investment) and a
distribution of installer income more realistic than the lump-sum rebate.

\item \textbf{No within-country energy-share heterogeneity.}
\citet{Bayer2026}'s central distributional mechanism is that poor households
have higher energy shares; we have homothetic energy demand conditional on
the durable type, so we cannot speak to their distributional results.

\item \textbf{No investment margin.} Both \citet{Bayer2026} and
\citet{ARS2024} have richer supply sides. Its absence here likely makes the
transfer look more expansionary, and the supply shock more inflationary, than
it would be with capital to smooth the adjustment.

\end{enumerate}

\appendix

\section{Summaries of the three reference papers}
\label{sec:summaries}

These summaries are in our own words, for internal reference. They are
deliberately compact---roughly one page each---and emphasise the features that
matter for our comparison in Section~\ref{sec:lit}.

\subsection{\citet{ARS2024}, ``Managing an Energy Shock''}

\citet{ARS2024} ask how an energy-importing economy should respond to an energy
price shock, and whether household heterogeneity changes the answer. The model
is a small open economy HANK built on the exchange-rate model of Auclert,
Rognlie, Souchier and Straub (2021), extended with an energy good. There are
three goods: a domestically produced home good, an imported foreign good, and
imported energy. Households face idiosyncratic productivity risk, cannot insure
it, and hold a single asset through a mutual fund; preferences are CRRA in a
CES consumption basket that nests energy against a home/foreign bundle, with a
low energy substitution elasticity ($\eta_E=0.1$) and a total substitution
index $\chi=0.3$. Crucially, all households share the \emph{same} basket and
the \emph{same} price index---the energy technology is fixed and homothetic.
Wages are sticky with a real-wage-stabilisation motive; the monetary baseline
is a rule that holds the real interest rate constant. Energy is supplied by
price-taking foreign firms with an intertemporal extraction margin, so in the
single-country equilibrium the world energy price is exogenous.

The analytical core is a decomposition of the output response into an
expenditure-switching channel (positive) and a real-income channel (negative)
scaled by the matrix of intertemporal MPCs. In the complete-markets
representative-agent benchmark only expenditure switching survives, so the
shock is expansionary; with realistic MPCs and low $\chi$ the real-income
channel dominates and the shock is contractionary---their central result, and
the one we inherit. A neutrality result holds at $\chi=1$, where the two
channels cancel. With real wage rigidity a wage--price spiral can appear, but
the real wage still falls.

On policy, ARS study monetary tightening and three fiscal tools---energy
subsidies, targeted transfers (indexed to counterfactual energy use), and
untargeted transfers, all deficit-financed and repaid by a labour tax. For a
country acting alone, subsidies are the most effective domestic stabiliser and
tame inflation (they cap the price that drives the real-wage loss); transfers
support demand but are more inflationary. The decisive twist is
\emph{cross-country}: monetary tightening has a positive externality (it lowers
world energy demand and price), while subsidies have a large negative
externality (coordinated subsidies make world energy demand nearly inelastic,
spike the world price, and become self-defeating). Their policy recommendation
is coordinated tightening plus transfers targeted to the poor, avoiding
subsidies. They also show the impulse responses are not strongly nonlinear at
empirically relevant shock sizes.

For our purposes the key features are: fixed homothetic energy technology, a
single energy price and CPI, the constant-real-rate monetary rule, and the fact
that the only ``cost'' of a subsidy in their world is a cross-country
externality---there is no transition margin.

\subsection{\citet{Bayer2026}, ``Hicks in HANK'' / ``Redistribution Within and
Across Borders''}

\citet{Bayer2026} evaluate the fiscal response to the European energy crisis in
a two-country HANK model of a monetary union, calibrated with Germany as
``Home'' (one-third of the union) and the rest of the area as ``Foreign''
(calibrated to Italy). The model is a medium-scale, two-asset HANK: households
hold a liquid union-wide bond and illiquid capital, face idiosyncratic income
risk, and have GHH preferences; there is an investment margin with capacity
utilisation. Energy enters both production (a CES bundle with the
capital--labour aggregate, elasticity 0.2) and household consumption (a CES
bundle with the physical good, elasticity 0.1). A distinctive and, for them,
central feature is within-country heterogeneity in the \emph{energy share}:
households are of high- or low-energy-intensity type, following a persistent
Markov process correlated with income, so that poorer households spend a larger
share on energy and suffer larger real-income losses. Monetary policy is a
standard union-wide Taylor rule; each national fiscal authority taxes labour
and profits and stabilises its debt.

The defining assumption is that the total union energy supply is
\emph{inelastic}: import capacity is fixed in the short run, so a fixed quantity
is set exogenously and a common price clears the market. The crisis is a 20\%
cut in supply for six quarters, in response to which the energy price rises
roughly fivefold, inflation rises about four points, and output falls around
one percent, with consumption falling most for poor, high-energy households.

They compare two policies, each implemented in Home only so that spillovers can
be measured: a subsidy that caps the retail energy price at its pre-crisis
level, and a Hicks/Slutsky transfer indexed to pre-crisis energy consumption
(to households and firms). The subsidy is the more effective domestic
stabiliser but is fiscally expensive and \emph{beggar-thy-neighbour}: by
holding Home's energy demand up it pushes the common price higher, amplifying
the crisis in Foreign. The transfer stabilises less at Home but generates
essentially no spillover and is cheaper. On welfare (consumption-equivalent
variation), the subsidy \emph{reduces} Home welfare---subsidising energy under a
supply contraction is inefficient because it suppresses both intratemporal and
intertemporal substitution---while the transfer \emph{raises} Home welfare, its
cost being mostly the distortion of the taxes that finance it. Their headline
contrast with ARS (subsidy more stabilising than transfer for output) they
attribute to the Taylor rule and the two-country structure, versus ARS's
constant-real-rate rule.

For our purposes the key features are: fixed union-wide quantity (our inelastic
closure), the two-country spillover mechanism that we cannot represent, the
within-country energy-share heterogeneity that we do not have, and the
substitution-preserving property of the transfer, which we share.

\subsection{\citet{Langot2026}, ``The Macroeconomic and Redistributive Effects
of the French Tariff Shield''}

\citet{Langot2026} evaluate the French \emph{bouclier tarifaire}---the 2022--23
cap on gas and electricity prices---in a HANK model of France as a single
economy within the euro area. The model extends Auclert, Rognlie, Souchier and
Straub with energy in both household consumption and production. Households face
idiosyncratic productivity risk and hold a single asset; preferences are CRRA
with a CES energy nest, and the energy substitution elasticity is $\eta_E=0.5$,
markedly higher than in ARS or Bayer. A distinctive feature is an
\emph{incompressible} (subsistence) level of energy consumption: only energy
above the subsistence level yields utility. This generates non-homotheticity
without non-homothetic preferences---poorer households devote a larger share to
incompressible energy and have a lower price elasticity of energy demand, which
is the source of the distributional results. The ECB Taylor rule is
down-weighted to reflect France's small weight in euro-area inflation. In the
policy simulations the energy price is taken as exogenous (from the
government's forecasts), so France is effectively a price taker facing
perfectly elastic supply.

The methodology is distinctive and worth flagging: rather than clean MIT-shock
impulse responses, the paper performs an \emph{ex-ante, conditional-forecast}
evaluation. Using the sequence-space Jacobian, the authors back out the
sequence of structural shocks (energy price, government spending, transfers,
plus latent demand, markup and fiscal shocks) that makes the model reproduce
the French government's official published forecasts for output, inflation and
the debt ratio through 2027. Policies are then evaluated as counterfactuals
that hold this shock sequence fixed and change one instrument.

The main finding is that the tariff shield supports growth (about 1.9\% average
over 2022--23 against 1.0\% without it) and curbs inflation (5.6\% against
7.2\%), at a fiscal cost near 2\% of GDP per year and a modest rise in the debt
ratio, while containing---though not eliminating---the rise in consumption
inequality even though it is untargeted. The mechanism behind its effectiveness
is a \emph{labour-cost} channel: by holding down the retail energy price the
shield prevents the price--wage spiral, contains the cost of labour, and
thereby supports employment and output. On this macro criterion the shield
beats both a faster wage indexation (more inflationary, worse for output) and a
targeted transfer to finance incompressible energy consumption; the targeted
transfer, however, reduces inequality by more.

For our purposes the key features are: exogenous/elastic energy price (our
price-taking closure), the labour-cost channel behind the pro-shield result
(which our wage block only partly captures), the higher intensive substitution
elasticity, and again a fixed energy technology with no transition margin.

\begin{thebibliography}{9}

\bibitem[Auclert et al.(2024)]{ARS2024}
Auclert, A., H. Monnery, M. Rognlie, and L. Straub (2024).
``Managing an Energy Shock: Fiscal and Monetary Policy.''
In \emph{Heterogeneity in Macroeconomics: Implications for Monetary Policy},
Series on Central Banking, Analysis, and Economic Policies, Vol. 30,
Central Bank of Chile, 39--108.

\bibitem[Bayer et al.(2026)]{Bayer2026}
Bayer, C., A. Kriwoluzky, G. J. M\"uller, and F. Seyrich (2026).
``Redistribution Within and Across Borders: The Fiscal Response to an Energy
Shock.'' \emph{Journal of Political Economy Macroeconomics}, forthcoming.
Earlier version: ``Hicks in HANK,'' CEPR DP 18557.

\bibitem[Iskhakov et al.(2017)]{DCEGM2017}
Iskhakov, F., T. H. J{\o}rgensen, J. Rust, and B. Schjerning (2017).
``The Endogenous Grid Method for Discrete-Continuous Dynamic Choice Models with
(or without) Taste Shocks.'' \emph{Quantitative Economics} 8(2), 317--365.

\bibitem[Langot et al.(2026)]{Langot2026}
Langot, F., S. Malmberg, F. Tripier, and J.-O. Hairault (2026).
``The Macroeconomic and Redistributive Effects of the French Tariff Shield.''
\emph{Journal of Political Economy Macroeconomics}, forthcoming.
Earlier version: CEPREMAP Working Paper 2305.

\end{thebibliography}

\end{document}

